% !TeX root = elegantnote-en.tex
\section{Aromatic Compounds}

\subsection{Introduction: The Discovery of Benzene 20260312 20260317}

\begin{enumerate}
    \item 
\end{enumerate}

\subsection{The Structure and Properties of Benzene 20260312 20260317}

\begin{enumerate}
    \item Kekulé structure, failed to explain existence of only one isomer of 1,2-dichlorobenzene.
    \item The Resonance Representation: a resonance hybrid between the two Kekulé structures, represented by drawing a circle inside the six-membered ring as a combined representation.
    \item Structure of Benzene: Each \textit{sp}$ ^2 $ hybridized C in the ring has an unhybridized \textit{p} orbital perpendicular to the ring which overlaps around the ring.
    \item An aromatic compound is a cyclic compound containing some number of conjugated double bonds and having an unusually large resonance energy.
    \item The Unusual Reactions of Benzene: will not undergo the typical reactions of polyenes.
    \item \textbf{Unusual Addition of Bromine to Benzene}
    \begin{enumerate}
        \item When bromine adds to benzene, a catalyst such as \textbf{\ch{FeBr3}} is needed.
        \item Substitution of a hydrogen by bromine.
        \item No addition of Br2 to the double bond.
    \end{enumerate}
    \item The difference between the predicted and the observed value  of hydrogenation is called the resonance energy.
    \item Annulenes
    \begin{enumerate}
        \item Annulenes are hydrocarbons with alternating single and double bonds.
        \item Cyclobutadiene is [4]annulene, Benzene is [6]annulene Ccyclooctatetraene is [8]annulene.
        \item Reactivity: Cyclobutadiene is so reactive that it dimerizes before it can be isolated. Cyclooctatetraene adds \ch{Br2} readily to the double bonds.
        \item tub conformation of cyclooctatetraene
    \end{enumerate}
\end{enumerate}

\subsection{The Molecular Orbitals of Benzene 20260317}

\begin{enumerate}
    \item Six overlapping p orbitals must form six molecular orbitals.
    
    Three will be bonding, three antibonding.
    
    Lowest energy MO will have all bonding interactions, no nodes.
    
    As energy of MO increases, the number of nodes increases.
    \item First MO of Benzene
    
    Intermediate MO of Benzene

    \item \textbf{Electronic Energy Diagram} for Benzene: Six electrons fill three bonding pi orbitals. All bonding orbitals are filled (“closed shell”), an extremely stable arrangement.
\end{enumerate}

\subsection{The Molecular Orbital Picture of Cyclobutadiene 20260317}

\begin{enumerate}
    \item Electronic Energy Diagram for Cyclobutadiene: Following Hund’s rule, two electrons are in separate nonbonding molecular orbitals. This diradical would be very reactive.
    \item The Polygon Rule: The molecular orbital energy diagram of a regular, completely conjugated cyclic system has the same polygonal shape as the compound, with one vertex (the all-bonding MO) at the bottom.
    \item However, cyclooctatetraene will not follow the Polygon Rule (planar) but will form a tub conformation to be more stable.
\end{enumerate}

\subsection{Aromatic, Antiaromatic, and Nonaromatic Compounds 20260317}

\begin{enumerate}
    \item Aromatic compounds are those that meet the following criteria:
    \begin{enumerate}
        \item Structure must be cyclic with conjugated pi bonds.
        \item Each atom in the ring must have an unhybridized \textit{p} orbital (\textit{sp}$ ^2 $ or \textit{sp}).
        \item The \textit{p} orbitals must overlap continuously \textbf{around the ring}. Structure must be planar (or close to planar for effective overlap to occur).
        \item Delocalization of the pi electrons over the ring must lower the electronic energy.
    \end{enumerate}
    \item An antiaromatic compound is one that meets the first three criteria, but delocalization of the pi electrons over the ring increases the electronic energy. Compounds or ions.
    \item Nonaromatic compounds do not have a continuous ring of overlapping \textit{p} orbitals and may be nonplanar.
\end{enumerate}

\subsection{Hückel’s Rule 20260317}

\begin{enumerate}
    \item To qualify as aromatic or antiaromatic, a cyclic compound must have a continuous ring of overlapping p orbitals, usually in a planar conformation.
    \item Once the criteria are met, Hückel’s rule applies:
    
    If the number of pi electrons in the cyclic system is: the system is aromatic. (4N) the system is antiaromatic. N is an integer, commonly 0, 1, 2, or 3.
    
    \item 4N Annulenes
    \begin{enumerate}
        \item {[4]}annulene is antiaromatic.
        \item {[8]}annulene would be antiaromatic, but it’s not planar, so it’s nonaromatic.
        \item Larger 4N annulenes are not antiaromatic because they are flexible enough to become nonplanar.
    \end{enumerate}
    
    \item 4N+2 Annulenes
    \begin{enumerate}
        \item {[6]}annulene (benzene) is aromatic.
        \item {[10]}annulene is aromatic except for the isomers that are nonplanar. (two hydrogen atoms interfere with each other)
        \item Some larger 4N+2 annulenes can achieve planar conformation, they are aromatic.
    \end{enumerate}
    
\end{enumerate}

\subsection{MO Derivation of Hückel’s Rule 20260317}

\begin{enumerate}
    \item Aromatic compounds have (4N+2) electrons and the orbitals are filled.
    \item Antiaromatic compounds have only 4N electrons and has unpaired electrons in two degenerate orbitals.
\end{enumerate}

\subsection{Aromatic ions 20260317 20260319}

\subsubsection{Cyclopentadienyl Ions}

\begin{enumerate}
    \item The cation has an empty p orbital, 4 electrons, so it is antiaromatic.
    \item The anion has a nonbonding pair of electrons in a \textit{p} orbital, 6 electrons, it is aromatic.
    \item Cyclopentadiene is acidic because deprotonation will convert it to an aromatic ion.
    \item Deprotonation of Cyclopentadiene: 
    
    By deprotonating the \textit{sp}$^3$ carbon of cyclopentadiene, the electrons in the p orbitals can be delocalized over all five carbon atoms and the compound would be aromatic.
    
    Cyclopentadiene is acidic because deprotonation will convert it to an aromatic ion.
    \item Resonance Forms of Cyclopentadienyl Ions
    
    The resonance picture gives \textbf{a misleading suggestion} of stability.
\end{enumerate}

\subsubsection{The Cycloheptatrienyl Ions}

\begin{enumerate}
    \item The cycloheptatrienyl cation has 6 pi electrons and an empty \textit{p} orbital. 
    \item The cycloheptatrienyl cation is easily formed by treating the corresponding alcohol with dilute (0.01 molar) aqueous sulfuric acid.
    \item The cycloheptatrienyl cation is called the tropylium ion. This aromatic ion is much less reactive than most carbocations.
\end{enumerate}

\subsubsection{The Cyclooctatetraene Dianion}

\begin{enumerate}
    \item Cyclooctatetraene reacts with potassium metal to form an aromatic dianion.
    \item The dianion has 10 pi electrons and is aromatic.
    \item Biphenyl has the following structure. p752
    \item Hexahelicene p754
\end{enumerate}

\begin{example} %16-12
    Explain why each compound or ion should be aromatic, antiaromatic, or nonaromatic.
    \begin{enumerate}
        \item the cyclononatetraene cation, anti
        \item the cyclononatetraene anion, aro
        \item the [16]annulene dianion, aro if planar
        \item the [18]annulene dianion, anti
        \item non
        \item the [20]annulene dication, aro if planar
    \end{enumerate}
\end{example}

\begin{example} %16-13
    The following hydrocarbon has an unusually large dipole moment. Explain how a large dipole moment might arise.

    \noindent \textbf{SOLUTION} \quad The reason for the dipole can be seen in a resonance form distributing the electrons to give each ring 6 pi electrons. This resonance picture gives one ring a negative charge and the other ring a positive charge.
\end{example}

\begin{example} %16-14
    When 3-chlorocyclopropene is treated with \ch{AgBF4}, \ch{AgCl} precipitates. The organic product can be obtained as a crystalline material, soluble in polar solvents such as nitromethane but insoluble in hexane. When the crystalline material is dissolved in nitromethane containing \ch{KCl}, the original 3-chlorocyclopropene is regenerated. Determine the structure of the crystalline material, and write equations for its formation and its reaction with chloride ion.
\end{example}

\begin{example} %16-15
    The polarization of a carbonyl group can be represented by a pair of resonance structures:
\end{example}

\subsubsection{Summary of Annulenes and Their Ions}

\begin{enumerate}
    \item The 2, 6, and 10 electron systems are aromatic, while the 4 and 8 electron systems are antiaromatic if they are planar.
\end{enumerate}

\subsection{Heterocyclic Aromatic Compounds 20260319}

\subsubsection{Pyridine}

\begin{enumerate}
    \item Pyridine has a nitrogen atom in place of one of the six \chemfig{C-H} units of benzene.
    \item Pyridine has a six-membered heterocyclic ring with six pi electrons.
    \item The nonbonding electrons are in an \textit{sp}$^2$ hybrid orbital in the plane of the ring, perpendicular to the pi system and do not overlap with it.
    \item Pyridine is basic because of its available pair of nonbonding electrons.
    \item The pyridinium ion is still aromatic because the additional proton has no effect on the electrons of the aromatic sextet.
\end{enumerate}

\subsubsection{Pyrrole}

\begin{enumerate}
    \item Pyrrole is an aromatic five-membered heterocycle, with one nitrogen atom and two double bonds.
    \item The pyrrole nitrogen atom is \textit{sp}$^2$ hybridized, and has a lone pair of electrons in the unhybridized \textit{p} orbital.
    \item The unhybridized orbital overlaps with the \textit{p} orbitals of the carbon atoms to form a continuous ring.
    \item Pyrrole is a much weaker base than pyridine. To form a bond to a proton requires the use of one of the electron pairs in the aromatic sextet.
    \item Basicity: Delocalized lone pair of electrons is not basic, usually with a hydrogen atom on the nitrogen atom.
\end{enumerate}

\subsubsection{Pyrimidine and Imidazole}

\begin{enumerate}
    \item Pyrimidine is a six-membered heterocycle with two nitrogen atoms situated in a 1,3- arrangement.
    \item Imidazole is an aromatic five-membered heterocycle with two nitrogen atoms.
    \item Purine has an imidazole ring fused to a pyrimidine ring. Purine has three basic nitrogen atoms and one pyrrole-like nitrogen.
\end{enumerate}

\begin{example} %16-18
    The proton NMR spectrum of 2-pyridone gives the chemical shifts shown.
\end{example}

\subsubsection{Furan and Thiophene 20260319}

\begin{enumerate}
    \item Furan is an aromatic five-membered heterocycle, and the heteroatom is oxygen.
    \item Thiophene is similar to furan, with a sulfur atom in place of the furan oxygen. The sulfur atom uses an unhybridized 3\textit{p} orbital to overlap with the 2\textit{p} orbitals on the carbon atoms.
\end{enumerate}

\begin{example} %16-19
    Explain why each compound is aromatic, antiaromatic, or nonaromatic.
\end{example}

\begin{example} %16-20 硼嗪
    Borazole, \ch{B3N3H6}, is an unusually stable cyclic compound. Propose a structure for borazole, and explain why it is aromatic.
\end{example}

%16-32, 34, 35

\subsection{Polynuclear Aromatic Hydrocarbons 20260324}

\begin{enumerate}
    \item Composed of two or more fused benzene rings, sharing two carbon atoms and the bond between them.
    \item Naphthalene: the simplest fused aromatic compound, consisting of two fused benzene rings.
    \item Anthracene and Phenanthrene: As the number of fused aromatic rings increases, the resonance energy per ring continues to decrease and the compounds become more reactive.
    \item Larger Polynuclear Aromatic Hydrocarbons: have more fused rings than anthracene and phenanthrene, and they have less resonance energy per ring and are more reactive.
\end{enumerate}

\subsection{Aromatic Allotropes of Carbon 20260324}

\subsubsection{Allotropes of Carbon: Diamond}

\begin{enumerate}
    \item Amorphous carbon refers to charcoal, soot, coal, and carbon black.
    \item Diamond has a crystalline structure containing tetrahedral carbon atoms linked together in a three-dimensional lattice.
\end{enumerate}

\subsubsection{Graphite}

\begin{enumerate}
    \item Graphite has the layered planar structure.
    \item Each layer of graphite is a nearly infinite lattice of fused aromatic rings. No bonds but van der Waals forces hold the layers together.
\end{enumerate}

\subsubsection{Fullerenes}

\subsection{Fused Heterocyclic Compounds 20260324}

Electrostatic potential map shows electron-rich (red) and electron-poor (blue).

\subsection{Nomenclature of Benzene Derivatives 20260324}

\begin{enumerate}
    \item The following compounds are usually called by their historical common names, and almost never by the systematic IUPAC names:
    
    phenol (benzenol), toluene (methylbenzene), aniline (benzenamine), anisole (methoxybenzene), styrene (vinylbenzene), acetophenone (methyl phenyl ketone), benzaldehyde, benzoic acid 

    \item Many compounds are named as derivatives of benzene.

    \textit{tert}-butylbenzene, nitrobenzene, ethynylbenzene (phenylacetylene), benzenesulfonic acid

    \item Disubstituted benzenes are named using the prefixes \textbf{ortho-} (1,2), \textbf{meta-} (1,3), and \textbf{para-} (1,4) to specify the substitution patterns. These terms are abbreviated \textit{o}-, \textit{m}-, and \textit{p}-.
    
    \item With three or more substituents on the benzene ring, numbers are used to indicate their positions, and give the lowest possible numbers to the substituents. The carbon atom bearing the functional group that defines the base name.
    
    \item Many disubstituted benzenes (and polysubstituted benzenes) have historical names. 
    
    \textit{m}-xylene,

    \item When the benzene ring is named as a substituent on another molecule, it is called a \textbf{phenyl group}. 
\end{enumerate}

\subsection{Physical Properties of Benzene and Its Derivatives 20260324}

\subsection{Spectroscopy of Aromatic Compounds 20260324 20260326}

\subsubsection{Infrared Spectroscopy}

\begin{enumerate}
    \item Aromatic compounds are readily identified by their infrared spectra because they show a characteristic \chemfig{C=C} stretch around \qty{1600}{\per\centi\meter}.
    \item s: strong, b: broad, m: medium
    \item 4000 - 2500 (\chemfig{X-H}, \chemfig{O-H}, \chemfig{N-H}, \chemfig{C-H}),
    
    2500 - 2000 (\chemfig{C~C}, \chemfig{C~N}),
    
    2000 - 1500 (\chemfig{C=C}, \chemfig{C=O}),
    
    1500 - 400 (fingerprint)
\end{enumerate}

\subsubsection{NMR Spectroscopy}

\subsubsection{Mass Spectrometry}

\begin{enumerate}
    \item The benzyl cation may rearrange to give the aromatic tropylium ion. 
    \item Alkylbenzenes frequently give ions corresponding to the tropylium ion at $ m/z $ 91.
\end{enumerate}

\subsubsection{Ultraviolet Spectroscopy}

\begin{enumerate}
    \item The most characteristic part of the spectrum is the band centered at \qty{254}{\nano\meter}, called the \textbf{benzenoid band}.
    \item About three to six small, sharp peaks (called fine structure) usually appear in this band.
\end{enumerate}
homework 16-24 27 28
